% Ders Kodu ve Adı (Dönem) [END 418 Stokastik Süreçler (2025 Bahar)]
% Bölüm [İMÜ MDBF Endüstri Mühendisliği Bölümü]

\documentclass[12 pt]{exam}

% ders, sınav, tarih, zaman, süre (dakika cinsinden)
\newcommand{\course}{Ders Kodu ve Adı} % [END 418 Stokastik Süreçler]
\newcommand{\exam}{SINAV ÇEŞİDİ \\ Klasik Kısım} % iki kısımdan oluşan sınavlarda klasik kısmın başlığı
%\newcommand{\exam}{SINAV ÇEŞİDİ \\ Bilgisayarlı Kısım} % iki kısımdan oluşan sınavlarda bilgisayarlı kısmın başlığı
%\newcommand{\exam}{\Large{SINAV ÇEŞİDİ}} % tek kısımdan oluşan sınavlarda başlık [ARA SINAV, FİNAL, BÜTÜNLEME, TEK DERS SINAVI]
%\newcommand{\exam}{MAZERET \\ ARA SINAVI} % tek kısımdan oluşan mazeret ara sınavı başlığı
\newcommand{\examdate}{Tarih} % [7 Mart 2025]
\newcommand{\examtime}{Saat} % [9.00] (\time komutunu yeniden tanımlayarak [yani \renewcommand{\time}{**.**} diyerek] devam edersek tikz paketini kullanırken hata veriyor, o yüzden \date ve \time komutlarını yeniden tanımlamak yerine yeni \examdate ve \examtime komutlarını kullanıyoruz)
\newcommand{\duration}{0} % [50]
\newcommand{\type}{Grup} % [A, B]

\printanswers % cevapları yazdırmak için (varsayılan yazdırmamak)

\usepackage[utf8]{inputenc}
\usepackage[turkish, shorthands = :!]{babel}
\usepackage[T1]{fontenc}
\usepackage{lmodern}

% sayfa kenar boşluklarını 1 cm'ye çekmek için (kenar boşluklarını geometry paketini kullanarak 1 cm yapacak olursak sayfa numaraları görünmüyor, görünmesi için boşluk biraz daha büyük olmalı)
\extrawidth{3.08 cm}
\extraheadheight{-1.54 cm}
\extrafootheight{-1.54 cm}

% ondalıkları virgül, binleri nokta ile ayırabilmek ve birimli sayıları düzgün bir biçimde yazabilmek için
\usepackage{siunitx}
\sisetup{output-decimal-marker = {,}}
\sisetup{group-separator = {.}}
\sisetup{group-minimum-digits = 4}

\usepackage{mathtools, lipsum}

\usepackage{graphicx}
\graphicspath{{./auxiliary/}}

% cevap anahtarının varsayılan doğru cevap gösterimi matematiksel cevapları kalınlaştırmadığı için doğru cevapları renklendirmek istersek
\usepackage[dvipsnames]{xcolor}
\CorrectChoiceEmphasis{\color{Green}}

% "points" yerine "puan" yazdırmak için
\pointname{\ puan}

% "solution" yerine "cevap" yazdırmak için
\renewcommand{\solutiontitle}{\noindent\textbf{Cevap:}\enspace}

% doğru-yanlış soruları için
\newcommand{\tf}[1][{}]{%
\fillin[#1][0.75 cm]%
}

% sayfa numaralarını ve grubu yazdırmak için
\pagestyle{foot}
\footer{}{\footnotesize\thepage/\numpages}{\footnotesize\type}

\begin{document}

\noindent
\begin{minipage}{0.35\textwidth}
\ifprintanswers
\itshape
\Large{Cevap Anahtarı}
\else
Ad Soyadı: \\ % sınav kâğıdında bölüm bilgisi istenmiyorsa (varsayılan)
Öğrenci Numarası: % sınav kâğıdında bölüm bilgisi istenmiyorsa (varsayılan)
%Öğrenci Numarası: \\ % sınav kâğıdında bölüm bilgisi isteniyorsa
%Bölüm: % sınav kâğıdında bölüm bilgisi isteniyorsa
\fi
\end{minipage}%
%
\begin{minipage}{0.30\textwidth}
\bfseries
\begin{center}
\exam
\end{center}
\end{minipage}%
%
\begin{minipage}{0.35\textwidth}
\begin{flushright}
\course \\ % sınav kâğıdında bölüm bilgisi istenmiyorsa (varsayılan)
%Dersin Kodu \\ % sınav kâğıdında bölüm bilgisi isteniyorsa
%Dersin Adı \\ % sınav kâğıdında bölüm bilgisi isteniyorsa
\examdate\ • \examtime\ • \qty{\duration}{dk}.
\end{flushright}
\end{minipage}

\smallskip

\hrule

\medskip

\begingroup
\noindent
\footnotesize
\textit{% klasik sınavla ilgili açıklamalar (ara sınav, mazeret ara sınavı)
Sınav çoktan seçmeli sorulardan oluşmaktadır. Her sorunun tek bir doğru cevabı vardır. Yanlış doğruyu götürmeyecektir. Bir soru için birden çok şık işaretlenmemelidir, aksi hâlde o soruya cevap verilmemiş sayılır. Sınava arkalı önlü bir A4 çalışma kâğıdı ile girilebilir, bunun dışında herhangi bir kaynaktan yararlanmak yasaktır. Bilimsel hesap makinesi kullanılabilir; cep telefonu, akıllı saat vb. cihazlar kullanılamaz.
}
%\textit{% klasik sınavla ilgili açıklamalar (final, bütünleme, tek ders)
%Sınav çoktan seçmeli sorulardan oluşmaktadır. Her sorunun tek bir doğru cevabı vardır. Yanlış doğruyu götürmeyecektir. Bir soru için birden çok şık işaretlenmemelidir, aksi hâlde o soruya cevap verilmemiş sayılır. Sınava arkalı önlü iki A4 çalışma kâğıdı ile girilebilir, bunun dışında herhangi bir kaynaktan yararlanmak yasaktır. Bilimsel hesap makinesi kullanılabilir; cep telefonu, akıllı saat vb. cihazlar kullanılamaz.
%}
%\textit{% bilgisayarlı sınavla ilgili açıklamalar
%Sınav çoktan seçmeli sorulardan oluşmaktadır. Her sorunun tek bir doğru cevabı vardır. Yanlış doğruyu götürmeyecektir. Bir soru için birden çok şık işaretlenmemelidir, aksi hâlde o soruya cevap verilmemiş sayılır. Bilgisayar sınav boyunca uçak modunda tutulmalıdır. Sınav öncesinde bilgisayara yüklenen her türlü kaynaktan yararlanmak serbesttir.
%}
\endgroup

\smallskip

\hrule

% ek açıklama (\begin{questions}'tan hemen sonra gelirse dosya dizilmiyor)
%\bigskip
%\smallskip
%\textit{Aşağıdaki sorular R programının \texttt{fpp3} kütüphanesi kullanılarak cevaplandırılacaktır.}

\begin{questions}

%%%%%%% SAYFA 1 %%%%%%%

% çoktan seçmeli sorular

\question[**] Çoktan seçmeli bir soru... % A
%\question[**] Çoktan seçmeli bir soru... % B

\begin{oneparchoices}
\choice birinci şık
\CorrectChoice ikinci şık % A
\choice üçüncü şık % B
\end{oneparchoices}

\vspace{\stretch{1}} % burada 1'den başka sayılar da kullanılabiliyor

\question[**] Çoktan seçmeli bir soru...

\begin{choices}
\choice birinci şık
\CorrectChoice ikinci şık
\choice üçüncü şık
\end{choices}

\vspace{\stretch{1}}

\question Çoktan seçmeli ve parçalı bir soru...

\begin{parts}

\part[**] Birinci parça...

\begin{oneparchoices}
\choice birinci şık
\CorrectChoice ikinci şık
\choice üçüncü şık
\end{oneparchoices}

\vspace{\stretch{1}}

\part[**] İkinci parça...

\begin{oneparchoices}
\choice birinci şık
\CorrectChoice ikinci şık
\choice üçüncü şık
\end{oneparchoices}

\vspace{\stretch{1}}

\end{parts}

\question[**] Çoktan seçmeli bir soru...

\begin{oneparcheckboxes}
\choice birinci şık
\CorrectChoice ikinci şık
\choice üçüncü şık
\end{oneparcheckboxes}

\vspace{\stretch{1}}

\question[**] Çoktan seçmeli bir soru...

\begin{checkboxes}
\choice birinci şık
\CorrectChoice ikinci şık
\choice üçüncü şık
\end{checkboxes}

\vspace{\stretch{1}}

\newpage

%%%%%%% SAYFA 2 %%%%%%%

\begingroup
%
\newcommand \lf {0.75} % sol panelin genişliğini belirleyen çarpan
\newcommand \gf {0.045} % orta boşluğun genişliğini belirleyen çarpan
\newcommand \rf {0.17} % sağ panelin genişliğini belirleyen çarpan
% çarpanların toplamı -deneme yanılmayla bulduk- 0.965 olmalı
%
% sol panel
%
\begin{minipage}[t]{\lf\linewidth}
\vspace{0 cm}
% sol panelin içi
\question[**] \lipsum[1][1-6] % A
%\question[**] \lipsum[2][1-6] % B
%
\end{minipage}% sol panelin sonu
%
% orta boşluk
%
\hspace*{\gf\linewidth}
%
% sağ panel
%
\begin{minipage}[t]{\rf\linewidth}
\vspace{0 cm}
% sağ panelin içi
\begin{center}
\includegraphics[scale = 0.1]{logo}
\end{center}
%
\end{minipage}% sağ panelin sonu
%
\endgroup

\medskip

\begin{oneparchoices}
\choice birinci şık
\CorrectChoice ikinci şık % A
\choice üçüncü şık % B
\choice dördüncü şık
\choice beşinci şık
\choice altıncı şık
\end{oneparchoices}

\vspace{\stretch{1}}

% klasik soru

\question[**] Klasik soru...

\begin{solution}
Sorunun cevabı...
\end{solution}

\vspace{\stretch{1}}

% doğru-yanlış sorusu

\question[**] Doğru-yanlış sorusu...
% Aşağıdaki ifadelerin doğru mu (D) yanlış mı (Y) olduklarını sol taraflarındaki boşluğa yazınız.

\begingroup
\renewcommand{\partlabel}{•}
\begin{parts}
\part \tf[D] İlk ifade...
\part \tf[Y] İkinci ifade...
\end{parts}
\endgroup

\vspace{\stretch{1}}

% boşluk doldurma sorusu

\question[10] Boşluk doldurma sorusu...
% Aşağıdaki cümlelerde verilen boşluğu uygun kelime ya da kelimelerle doldurunuz.

\begingroup
\renewcommand{\partlabel}{•}
\begin{parts}
\part İlk \fillin[boşluklu] ifade...
\part İkinci \fillin[boşluk doldurmalı][4.25 cm] ifade...
\end{parts}
\endgroup

\vspace{\stretch{1}}

\end{questions}

\end{document}